
\frontslide{Supports de stockage}


\begin{frame}[fragile]{Techniques de stockage}

Contenu de ce cours\,:

\begin{redbullet}
   \item Où sont les données? Mémoires RAM, mémoire persistante.


  \item Les performances 

  \item Disques, SSD
\end{redbullet}

\vfill

\begin{beamerboxesrounded}{Vocabulaire}
Une "donnée" désigne une unité physique d'information: pour nous, une ligne d'une table,
codée sous forme d'un enregistrement (\textit{record}).
\end{beamerboxesrounded}

\vfill


\begin{blocgauche}
\alert{Ces diapositives correspondent au support en ligne disponible sur
le site \url{http://sys.bdpedia.fr/stock.html\#s1-supports-de-stockage}}
\end{blocgauche}

\end{frame}


%------------------------------------------------------------
%\subsection{Supports de stockage: les mémoires}
%------------------------------------------------------------

\begin{frame}{Problématique du stockage}

Un SGBD gère deux types de mémoire.

\begin{redbullet}
  \item La mémoire centrale, \alert{volatile}, rapide mais risque de perte.
  \item La mémoire secondaire, \alert{persistante}, stable mais lente.
\end{redbullet}

Les données sont \alert{toujours} en mémoire secondaire dans
des fichiers, et \alert{partiellement} copiées en mémoire centrale.

\vfill

\begin{blocgauche}
Une donnée, pour être traitée, doit impérativement être en mémoire centrale.

\medskip

Comment gérer au mieux ces deux ressources?
\end{blocgauche}
\end{frame}


\begin{frame}{Les mémoires d'un ordinateur}

Les mémoires dans un ordinateur forment une hiérarchie.

\vfill

\alert{Constat}: plus une mémoire est rapide, plus elle est chère , et moins elle est volumineuse.



\vfill

\figSlideWithSize{hiermem}{6cm}

\end{frame}


\begin{frame}{Performances des mémoires}

Critères essentiels:

\begin{redbullet}
  \item \alert{Temps d'accès}: connaissant  l'adresse d'une donnée, quel est le temps
    nécessaire pour aller la chercher  à l'emplacement mémoire indiqué par cette adresse? 
    
    Critère essentiel pour les opérations dites \alert{d'accès direct}.
    
  \item \alert{Débit}: volume de données lues par unité de temps
    dans le meilleur des cas.
    
     Critère essentiel pour les opérations dites de \alert{lecture séquentielle}.
    
\end{redbullet}

\vfill
\begin{blocgauche}
\alert{Constat}: tout accès aux données s'effectue avec l'une de ces deux opérations.
\end{blocgauche}
\end{frame} 
    
\begin{frame}
{Quelques ordres de grandeur}

Rappel: 1 Go = 1\,000 Mo  et 1 To = 1\,000 Go.

\vfill

\begin{center}
\begin{small}
\begin{tabular}{l|l|l|l}
{\bf Mémoire} & {\bf Taille} & {\bf Temps d'accès} & {\bf Débit}\\
\hline
{\it cache} & Qq Mo &  $\approx 10^{-8}$ (10 nanosec.) & 10+ Go/s\\
Principale (RAM)  &  Qq Go &  $\approx 10^{-8} - 10^{-7}$ (10-100 nanosec.) & Qq Go/s \\
Disque  & Qq To &  $\approx 10^{-2}$ (10 millisec.) & 100 Mo/s\\
SSD & Qq To &  $\approx 10^{-4}$ (0,1 millisec.) & 1+ Go/s\\
\hline
\end{tabular}
\end{small}
\end{center}

\vfill
\begin{blocgauche}
Un accès \alert{direct} disque est (environ) un million de fois
plus coûteux qu'en mémoire principale\,!

\smallskip

Le débit d'un disque est 20, 30, 40 fois plus lent que le débit RAM.

\smallskip

Un SSD est 10 à  100 fois plus performant qu'un disque magnétique.
\end{blocgauche}

\end{frame}



\begin{frame}{Importance pour les bases de données}

 Un SGBD doit ranger sur disque les données\,;
\begin{redbullet}
   \item parce qu'elle sont trop volumineuses (de moins en moins vrai)\,;
   \item parce qu'elles sont \emph{persistantes} (et doivent survivre
      à un arrêt du système).
\end{redbullet}

\vfill
Le SGBD doit \alert{toujours} amener 
      les données en mémoire pour les traiter.

\vfill

Si possible, les \alert{données utiles} devraient résider
         le plus possible en mémoire.

\vfill

\begin{blocgauche}
\alert{Principe général}:
la capacité à placer les données \alert{utiles} au plus près du processeur (donc en RAM si possible)
est un facteur essentiel de performance.
\end{blocgauche}

\end{frame}

%\subsection{Les disques magnétiques}


\begin{frame}
{Organisation d'un disque}

\alert{Disque}\,: surface magnétique, stockant des 0 ou des 1,
          divisé en secteurs.

\vfill

%Par extension\,: ensemble de surfaces (simple ou double-face)
%                     ayant mêmes caractéristiques.

\alert{Dispositif}\,: les surfaces sont entraînées dans un mouvement de rotation\,;
              les têtes de lecture se déplacent dans un plan fixe.

\begin{blocgauche}
\begin{redbullet}
  \item Le \alert{bloc} est un ensemble de secteurs, sa
              taille est en général un multiple de 512\,;

  \item La \alert{piste} est l'ensemble des blocs d'une surface
       lus au cours d'une rotation\,;

  \item le  \alert{cylindre} est un ensemble de pistes situées
                    sous les têtes de lecture.
\end{redbullet}
\end{blocgauche}
\end{frame}

\begin{frame}
{Structure d'un disque}

\figSlide{disques}

\end{frame}

\begin{frame}
\frametitle{Disque = accès \alert{semi-direct} = latence}

Adresse = numéro du disque\,;
 de la piste où se trouve le bloc\,;
 du numéro du bloc sur la piste.

\vfill

\begin{redbullet}
  \item \alert{Délai de positionnement} pour placer
              la tête sur la bonne piste\,:

  \item \alert{Délai de latence} pour attendre que le bloc
           passe sous la tête de lecture\,;

  \item \alert{Temps de transfert} pour attendre que le (ou les)
                bloc(s) soient lus et transférés.
\end{redbullet}

\vfill

\begin{blocgauche}
\alert{La latence d'un disque}:
on ne peut pas lire une donnée avant qu'elle passe
sous une tête de lecture: la \alert{latence} (qq ms) 
est un facteur pénalisant.
\end{blocgauche}


\end{frame}

\begin{frame}{Granularité d'entrée/sortie: le bloc}

\alert{Granularité}: on lit toujours au moins un bloc,
    même si on ne veut qu'un octet\,; le \alert{bloc} est \alert{l'unité d'entrée/sortie}\,! 

\medskip

  $\Rightarrow$ le remplissage des blocs, si possible  avec des données
     ``proches'' de la données requise initialement, est essentiel.
     
\vfill

\alert{Localité}\,: la latence en pratique dépend de la proximité
   des adresses successives demandées par l'application. 

\begin{blocgauche}
\alert{Principe}: si deux données sont \alert{proches} du point de vue applicatif,
  alors elles doivent être proches sur le disque (idéalement, dans
   le même bloc).
\end{blocgauche}

\end{frame}


\begin{frame}{Résumé}

La gestion des disques n'est pas directement sous le contrôle de l'administrateur. 

\vfill 
\alert{Retenir}:

\begin{redbullet}
  \item Accès (aléatoire) au disque\,: très lent par rapport à un accès en RAM (qq. ms).
  \item Débit du disque\,: faible (100 MO/s au mieux)
  \item SSD: 100 fois plus rapide qu'un disque (latence), 10 fois plus en débit
\end{redbullet}

\vfill

\begin{blocgauche}

\alert{Règle générale}: définir des données compactes (types) et stockées contiguement.

\medskip

Organisation des données, les algorithmes d'accès, la concurrence: tout est conçu pour minimiser
les accès aux disques, et surtout les accès \alert{aléatoires}.
\end{blocgauche}

\end{frame}
